% DeePC presentation content for IguanaTex / LaTeX slide equations
% Source files: deepc_data_collection.m and deepc_reg.m
%
% Suggested presentation flow:
% 1. Why DeePC for cart-pole?
% 2. Offline data collection
% 3. PRBS + stabilizing feedback
% 4. Hankel data matrices
% 5. Persistency of excitation
% 6. Regularized DeePC optimization
% 7. Receding-horizon implementation
% 8. Diagnostics and takeaways

% ------------------------------------------------------------
% Common IguanaTex preamble
% ------------------------------------------------------------
\usepackage{amsmath,amssymb,bm}
\newcommand{\R}{\mathbf{R}}
\newcommand{\Tini}{T_{\mathrm{ini}}}
\newcommand{\Npred}{N}
\newcommand{\Ndata}{N_{\mathrm{data}}}
\newcommand{\Ts}{T_s}
\newcommand{\Hankel}{\mathcal{H}}
\newcommand{\col}{\operatorname{col}}
\newcommand{\rank}{\operatorname{rank}}

% ============================================================
% Slide 1: DeePC story
% ============================================================
% Slide title:
% Data-Enabled Predictive Control for the Cart-Pole
%
% Speaker storyline:
% DeePC replaces explicit prediction by model matrices with prediction
% from one sufficiently rich input-output-state trajectory.

\[
\text{Model-based MPC:}\qquad
x_{k+1}=A_d x_k+B_d u_k
\]

\[
\text{DeePC:}\qquad
\text{future trajectory}
=
\text{linear combination of measured trajectories}
\]

% ============================================================
% Slide 2: Linearized cart-pole model used for simulation
% ============================================================
% Slide title:
% Plant Model Used to Generate and Validate Data

\[
x=
\begin{bmatrix}
p & \dot p & \phi & \dot \phi
\end{bmatrix}^{\top},
\qquad
y=
\begin{bmatrix}
p\\ \phi
\end{bmatrix}
\]

\[
\dot{x}(t)=Ax(t)+Bu(t),
\qquad
y(t)=Cx(t)+Du(t)
\]

\[
C=
\begin{bmatrix}
1&0&0&0\\
0&0&1&0
\end{bmatrix},
\qquad
D=
\begin{bmatrix}
0\\0
\end{bmatrix}
\]

\[
x_{k+1}=A_d x_k+B_d u_k,
\qquad
y_k=C_d x_k+D_d u_k,
\qquad
T_s=0.01~\mathrm{s}
\]

% ============================================================
% Slide 3: Why data collection needs feedback
% ============================================================
% Slide title:
% The Data Must Be Rich, But the Plant Is Unstable
%
% Main point:
% Strict open-loop PRBS can excite the system, but the inverted pendulum
% can leave the local linearized region. The script therefore collects data
% with stabilizing LQR feedback plus PRBS excitation.

\[
u_k = u_{\mathrm{fb},k}+u_{\mathrm{PRBS},k}
\]

\[
u_{\mathrm{fb},k}=-K_{\mathrm{data}}x_k
\]

\[
\boxed{
u_k=-K_{\mathrm{data}}x_k+u_{\mathrm{PRBS},k}
}
\]

\[
u_{\mathrm{PRBS},k}
\in
\left[-4,\,4\right],
\qquad
u_k \in \left[-10,\,10\right]\ \text{checked diagnostically}
\]

% ============================================================
% Slide 4: Offline data collected by deepc_data_collection.m
% ============================================================
% Slide title:
% One Continuous Trajectory Becomes the DeePC Dataset

\[
N_{\mathrm{data}}=3000,
\qquad
T_{\mathrm{ini}}=8,
\qquad
N=40,
\qquad
L=T_{\mathrm{ini}}+N=48
\]

\[
U \in \mathbf{R}^{1\times N_{\mathrm{data}}},
\qquad
Y \in \mathbf{R}^{2\times N_{\mathrm{data}}},
\qquad
X \in \mathbf{R}^{4\times(N_{\mathrm{data}}+1)}
\]

\[
Y=
\begin{bmatrix}
\text{cart position}\\
\text{pole angle}
\end{bmatrix},
\qquad
X=
\begin{bmatrix}
p\\ \dot p\\ \phi\\ \dot \phi
\end{bmatrix}
\]

% ============================================================
% Slide 5: Block Hankel construction
% ============================================================
% Slide title:
% Hankel Matrices Turn Data into Trajectory Libraries

\[
\mathcal{H}_L(w)=
\begin{bmatrix}
w_1 & w_2 & \cdots & w_{T-L+1}\\
w_2 & w_3 & \cdots & w_{T-L+2}\\
\vdots & \vdots & \ddots & \vdots\\
w_L & w_{L+1} & \cdots & w_T
\end{bmatrix}
\]

\[
\begin{bmatrix}
U_p\\U_f
\end{bmatrix}
=\mathcal{H}_L(U),
\qquad
\begin{bmatrix}
Y_p\\Y_f
\end{bmatrix}
=\mathcal{H}_L(Y),
\qquad
\begin{bmatrix}
X_p\\X_f
\end{bmatrix}
=\mathcal{H}_L(X)
\]

\[
U_p\in\mathbf{R}^{n_uT_{\mathrm{ini}}\times n_c},
\quad
U_f\in\mathbf{R}^{n_uN\times n_c},
\quad
X_p\in\mathbf{R}^{n_xT_{\mathrm{ini}}\times n_c},
\quad
X_f\in\mathbf{R}^{n_xN\times n_c}
\]

% ============================================================
% Slide 6: Past and future split
% ============================================================
% Slide title:
% The Same Coefficient Vector Must Explain Past and Future

\[
g\in\mathbf{R}^{n_c}
\]

\[
\underbrace{
\begin{bmatrix}
U_p\\X_p
\end{bmatrix}g
}_{\text{match measured past}}
=
\underbrace{
\begin{bmatrix}
u_{\mathrm{ini}}\\x_{\mathrm{ini}}
\end{bmatrix}
}_{\text{latest window}}
\]

\[
\underbrace{
\begin{bmatrix}
U_f\\X_f
\end{bmatrix}g
}_{\text{predicted future}}
=
\underbrace{
\begin{bmatrix}
u\\x
\end{bmatrix}
}_{\text{optimization variables}}
\]

% ============================================================
% Slide 7: Persistency of excitation
% ============================================================
% Slide title:
% Richness Check: Persistency of Excitation

\[
L_{\mathrm{PE}}=T_{\mathrm{ini}}+N+n_x
\]

\[
L_{\mathrm{PE}}=8+40+4=52
\]

\[
\operatorname{rank}\!\left(\mathcal{H}_{L_{\mathrm{PE}}}(U)\right)
=
n_u L_{\mathrm{PE}}
\]

\[
\operatorname{rank}\!\left(\mathcal{H}_{52}(U)\right)=52
\qquad
\text{is the target for } n_u=1.
\]

% ============================================================
% Slide 8: What is saved after data collection
% ============================================================
% Slide title:
% Offline Output Passed to the DeePC Controller

\[
\texttt{deepc\_cartpole\_dataset.mat}
\quad\Longrightarrow\quad
\{A_d,B_d,C_d,D_d,T_s,U,Y,X,U_p,U_f,X_p,X_f\}
\]

\[
\text{The controller does not use only }u_{\mathrm{PRBS}}.
\]

\[
\boxed{
\text{The Hankel matrices are built from the actual applied input }U.
}
\]

% ============================================================
% Slide 9: Online DeePC variables from deepc_reg.m
% ============================================================
% Slide title:
% Online Regularized DeePC Decision Variables

\[
g\in\mathbf{R}^{n_c},
\qquad
u\in\mathbf{R}^{n_uN},
\qquad
x\in\mathbf{R}^{n_xN},
\qquad
\sigma\in\mathbf{R}^{n_xT_{\mathrm{ini}}}
\]

\[
u_{\mathrm{ini}}\in\mathbf{R}^{n_uT_{\mathrm{ini}}},
\qquad
x_{\mathrm{ini}}\in\mathbf{R}^{n_xT_{\mathrm{ini}}},
\qquad
x_k\in\mathbf{R}^{n_x}
\]

% ============================================================
% Slide 10: DeePC behavioral constraints
% ============================================================
% Slide title:
% Behavioral Constraints Used in the QP

\[
U_p g = u_{\mathrm{ini}}
\]

\[
X_p g = x_{\mathrm{ini}}+\sigma
\]

\[
U_f g = u,
\qquad
X_f g = x
\]

\[
x_{0|k}=x_k
\]

% ============================================================
% Slide 11: Regularized DeePC objective
% ============================================================
% Slide title:
% Regularized DeePC Objective

\[
\bar Q=I_N\otimes Q_x,
\qquad
\bar R=I_N\otimes R_u
\]

\[
\min_{g,u,x,\sigma}
\ (x-r)^{\top}\bar Q(x-r)
+u^{\top}\bar R u
+\lambda_g\|g\|_1
+\lambda_{\sigma}\|\sigma\|_2^2
\]

\[
Q_x=\operatorname{diag}(200,\ 10,\ 1000,\ 100),
\qquad
R_u=0.05
\]

\[
\lambda_g=10,
\qquad
\lambda_{\sigma}=10^8
\]

% ============================================================
% Slide 12: Full QP formulation
% ============================================================
% Slide title:
% Optimization Solved at Every Control Step

\[
\begin{aligned}
\min_{g,u,x,\sigma}\quad
&(x-r)^{\top}\bar Q(x-r)
+u^{\top}\bar R u
+\lambda_g\|g\|_1
+\lambda_{\sigma}\|\sigma\|_2^2\\
\mathrm{s.t.}\quad
&U_p g=u_{\mathrm{ini}},\\
&X_p g=x_{\mathrm{ini}}+\sigma,\\
&U_f g=u,\\
&X_f g=x,\\
&x_{0|k}=x_k,\\
&u_{\min}\le u\le u_{\max},\\
&p_{\min}\le p\le p_{\max},\\
&\phi_{\min}\le \phi\le \phi_{\max},\\
&\phi_{N|k}=0.
\end{aligned}
\]

% ============================================================
% Slide 13: Constraints used in deepc_reg.m
% ============================================================
% Slide title:
% Constraint Values in the Implementation

\[
-9.5\le u_k\le 9.5
\]

\[
-0.40\le p_k\le 0.40
\]

\[
-0.15\le \phi_k\le 0.15
\]

\[
\phi_{N|k}=0
\]

% ============================================================
% Slide 14: Receding horizon loop
% ============================================================
% Slide title:
% DeePC Runs as a Receding-Horizon Controller

\[
\{u^\star,x^\star,g^\star,\sigma^\star\}
=
\operatorname{DeePC}
\left(
u_{\mathrm{ini}},
x_{\mathrm{ini}},
x_k,
r
\right)
\]

\[
u_k=u^\star_{0|k}
\]

\[
x_{k+1}=A_d x_k+B_d u_k
\]

\[
u_{\mathrm{ini}}
\leftarrow
\operatorname{shift}
\left(u_{\mathrm{ini}},u_k\right),
\qquad
x_{\mathrm{ini}}
\leftarrow
\operatorname{shift}
\left(x_{\mathrm{ini}},x_k\right)
\]

% ============================================================
% Slide 15: Closed-loop diagnostics
% ============================================================
% Slide title:
% What the Simulation Checks

\[
N_{\mathrm{sim}}=300
\]

\[
J_{\mathrm{DeePC}}
=
\sum_{k=0}^{N_{\mathrm{sim}}-1}
\left[
(x_k-r)^{\top}Q_x(x_k-r)
+u_k^{\top}R_u u_k
\right]
\]

\[
\text{reported diagnostics:}\quad
\begin{cases}
\text{successful QP solves},\\
\text{input/state constraint violations},\\
\|g\|_2,\ \|\sigma\|_2,\\
\text{average and worst solve time}.
\end{cases}
\]

% ============================================================
% Slide 16: Computation-time check
% ============================================================
% Slide title:
% Real-Time Feasibility Check

\[
T_s=0.01~\mathrm{s}=10~\mathrm{ms}
\]

\[
\text{real-time feasible if}\qquad
\max_k\ T_{\mathrm{step}}(k)<T_s
\]

\[
T_{\mathrm{step}}(k)
=
T_{\mathrm{solve}}(k)+T_{\mathrm{update}}(k)
\]

% ============================================================
% Slide 17: Final takeaway
% ============================================================
% Slide title:
% DeePC Pipeline: From Exciting Data to Constrained Control

\[
\boxed{
\text{PRBS + LQR data}
\ \Longrightarrow\
\text{Hankel trajectory library}
\ \Longrightarrow\
\text{regularized constrained QP}
\ \Longrightarrow\
\text{receding-horizon control}
}
\]

\[
\text{The key design choice is to preserve excitation while keeping the unstable plant near the linear region.}
\]
